% =====================================================================
%  IEL04 - Circuitos Eléctricos I
%  Semana 11: El circuito RLC en serie: reactancias opuestas, ley de voltajes de Kirchhoff y resonancia en serie
%  Institución Universitaria de Barranquilla (IUB)
%  Generada por src/presentaciones.py (diseño IUB 5). Compilador: pdfLaTeX (dos pasadas)
% =====================================================================
\documentclass[aspectratio=169,11pt,xcolor={table}]{beamer}

% ---------------------------------------------------------------------
%  Codificación, idioma y tipografía
% ---------------------------------------------------------------------
\usepackage[utf8]{inputenc}
\usepackage[T1]{fontenc}
\usepackage[spanish,es-nodecimaldot,es-noshorthands]{babel}
\usepackage{avant}                       % Sans geométrica (emula Avenir Next)
\renewcommand{\familydefault}{\sfdefault}
\renewcommand{\ttdefault}{pcr}           % Monoespaciada para código
\usepackage{textcomp}

% ---------------------------------------------------------------------
%  Paquetes
% ---------------------------------------------------------------------
\usepackage{amsmath,amssymb,mathtools}
\usepackage{siunitx}
\sisetup{output-decimal-marker={.}, per-mode=symbol}
\usepackage{booktabs,tabularx,multirow,array}
\usepackage{adjustbox}
\usepackage{tikz}
\usetikzlibrary{arrows.meta,positioning,calc,shapes.geometric,shapes.misc,
  decorations.pathreplacing,fit,backgrounds,babel}
\usepackage[siunitx,RPvoltages]{circuitikz}
\usepackage{pgfplots}
\pgfplotsset{compat=1.18}
\usepackage[most]{tcolorbox}
\usepackage{listings}

% ---------------------------------------------------------------------
%  Paleta institucional IUB (Manual de Identidad Visual 2025)
% ---------------------------------------------------------------------
\definecolor{iubwhite}{HTML}{FFFFFF}
\definecolor{iubnavy}{HTML}{121023}
\definecolor{iubyellow}{HTML}{FFDF2D}
\definecolor{iubblue}{HTML}{00ADE7}
\definecolor{iuborange}{HTML}{E39037}
\definecolor{iubgreen}{HTML}{3B9C6D}

% ---------------------------------------------------------------------
%  Configuración de Beamer
% ---------------------------------------------------------------------
\setbeamertemplate{navigation symbols}{}
\setbeamersize{text margin left=0.6cm, text margin right=0.6cm}
\setbeamercolor{normal text}{fg=iubnavy, bg=iubwhite}
\setbeamercolor{background canvas}{bg=iubwhite}
\setbeamercolor{structure}{fg=iubnavy}
\setbeamercolor{alerted text}{fg=iuborange}
\setbeamertemplate{itemize item}{\textcolor{iubblue}{\small$\blacktriangleright$}}
\setbeamertemplate{itemize subitem}{\textcolor{iuborange}{\scriptsize$\bullet$}}
\setbeamertemplate{enumerate item}{\textcolor{iubblue}{\bfseries\insertenumlabel.}}
\setbeamertemplate{frametitle}{}         % el título vive en el headline

% Parches: el headline se pre-mide antes del primer frame
\makeatletter
\providecommand{\insertframetitle}{}
\providecommand{\insertframesubtitle}{}
\makeatother

% ---------- Encabezado ----------
\setbeamertemplate{headline}{%
  \begin{tikzpicture}
    \useasboundingbox (0,0) rectangle (\paperwidth,1.05cm);
    \fill[iubnavy] (0,0) rectangle (\paperwidth,1.05cm);
    \fill[iubyellow] (0,0) rectangle (\paperwidth,0.05cm);
    \node[anchor=west, text=iubyellow, font=\large\bfseries]
      at (0.6cm,0.55cm) {\strut\insertframetitle};
    \node[anchor=east, text=iubyellow, font=\footnotesize\bfseries]
      at ([xshift=-0.6cm]\paperwidth,0.55cm) {IUB};
  \end{tikzpicture}}

% ---------- Pie de página ----------
\setbeamertemplate{footline}{%
  \begin{tikzpicture}
    \useasboundingbox (0,0) rectangle (\paperwidth,0.55cm);
    \fill[iubnavy] (0,0) rectangle (\paperwidth,0.55cm);
    \node[anchor=west, text=iubyellow, font=\tiny]
      at (0.6cm,0.275cm) {IEL04 \textperiodcentered{} Circuitos Eléctricos I};
    \node[anchor=center, text=iubyellow, font=\tiny]
      at (0.66\paperwidth,0.275cm) {Ing. Sergio Martinez Castilla};
    \node[anchor=east, text=iubyellow, font=\tiny\bfseries]
      at ([xshift=-0.6cm]\paperwidth,0.275cm)
      {Semana 11 \quad \insertframenumber\,/\,\inserttotalframenumber};
  \end{tikzpicture}}

% ---------------------------------------------------------------------
%  Cajas tcolorbox (texto blanco sobre la paleta complementaria)
%  \begin{caja}[opciones]{color}{título} ... \end{caja}
%  \begin{cajaplana}[opciones]{color} ... \end{cajaplana}
%  \begin{cardB|cardO|cardG|cardN}[opciones]{título} ... \end{cardX}
% ---------------------------------------------------------------------
\tcbset{
  iubbase/.style={enhanced, arc=1.5mm, boxrule=0pt, coltext=white, coltitle=white,
    fonttitle=\bfseries\footnotesize, fontupper=\footnotesize,
    left=1.8mm, right=1.8mm, top=1mm, bottom=1.2mm,
    toptitle=0.6mm, bottomtitle=0.6mm, halign=left,
    before upper={\hyphenpenalty=5000\setbeamercolor{itemize item}{fg=white}%
                  \setbeamercolor{itemize/enumerate body}{fg=white}%
                  \setbeamercolor{itemize/enumerate subbody}{fg=white}%
                  \setbeamercolor{itemize subitem}{fg=white}%
                  \setbeamercolor{enumerate item}{fg=white}%
                  \setbeamertemplate{itemize item}{\textcolor{white}{\small$\blacktriangleright$}}}}
}
\newtcolorbox{caja}[3][]{iubbase, colback=#2, colframe=#2,
  colbacktitle=#2!78!iubnavy, title={#3}, #1}
\newtcolorbox{cajaplana}[2][]{iubbase, colback=#2, colframe=#2, #1}
\newtcolorbox{cardB}[2][]{iubbase, colback=iubblue,   colframe=iubblue,
  colbacktitle=iubblue!78!iubnavy,   title={#2}, #1}
\newtcolorbox{cardO}[2][]{iubbase, colback=iuborange, colframe=iuborange,
  colbacktitle=iuborange!78!iubnavy, title={#2}, #1}
\newtcolorbox{cardG}[2][]{iubbase, colback=iubgreen,  colframe=iubgreen,
  colbacktitle=iubgreen!78!iubnavy,  title={#2}, #1}
\newtcolorbox{cardN}[2][]{iubbase, colback=iubnavy,   colframe=iubnavy,
  colbacktitle=iubnavy, coltitle=iubyellow, title={#2}, #1}

% Celda de encabezado de tabla (texto amarillo sobre \rowcolor{iubnavy}) y código en línea
\newcommand{\hd}[1]{\textcolor{iubyellow}{\bfseries #1}}
\newcommand{\thc}[1]{\textcolor{iubyellow}{\textbf{#1}}}
\newcommand{\cod}[1]{\texttt{\textbf{#1}}}

% ---------------------------------------------------------------------
%  Código sobre fondo oscuro:
%  \begin{codebox}[listing options app={language=Python,firstnumber=1}]{título}
% ---------------------------------------------------------------------
\lstdefinestyle{iubcode}{
  basicstyle=\ttfamily\fontsize{7}{8.4}\selectfont\color{white},
  keywordstyle=\color{iubblue}\bfseries,
  commentstyle=\color{iubgreen!55!white}\itshape,
  stringstyle=\color{iuborange!80!white},
  numbers=left, numberstyle=\tiny\color{iubyellow}, numbersep=6pt,
  showstringspaces=false, breaklines=true, tabsize=4,
  columns=fullflexible, keepspaces=true, upquote=true}
\newtcblisting{codebox}[2][]{enhanced, listing only, colback=iubnavy,
  colframe=iubnavy, colbacktitle=iubnavy!85!white, coltitle=iubyellow,
  fonttitle=\bfseries\scriptsize, title={#2}, arc=1.5mm, boxrule=0pt,
  left=6mm, right=1mm, top=0.5mm, bottom=0.5mm, toptitle=0.5mm, bottomtitle=0.5mm,
  listing options={style=iubcode}, #1}

% ---------------------------------------------------------------------
%  Estilos TikZ
% ---------------------------------------------------------------------
\tikzset{
  bloque/.style={rectangle, rounded corners=2mm, fill=#1, text=white,
    font=\scriptsize\bfseries, align=center, minimum height=1.1cm,
    text width=1.8cm, inner sep=1.2mm},
  flecha/.style={-{Stealth[length=2.2mm]}, line width=1pt, draw=iubnavy},
  arr/.style={flecha},
  term/.style={rectangle, draw=#1, line width=1.2pt, fill=white,
    text=iubnavy, font=\tiny\bfseries, minimum width=1.05cm,
    minimum height=0.5cm, inner sep=1pt},
  nodo/.style={rectangle, rounded corners=1mm, fill=#1, text=white,
    font=\scriptsize\bfseries, minimum size=0.75cm, align=center},
  cable/.style={line width=1.4pt, draw=#1},
  box/.style={rectangle, rounded corners=1.5mm, draw=none, text=white,
    font=\scriptsize\bfseries, align=center, minimum height=0.8cm, inner sep=3pt},
  bB/.style={box, fill=iubblue}, bO/.style={box, fill=iuborange},
  bG/.style={box, fill=iubgreen}, bN/.style={box, fill=iubnavy},
  dec/.style={diamond, aspect=1.9, fill=iuborange, text=white,
    font=\scriptsize\bfseries, align=center, inner sep=1pt},
  tag/.style={fill=#1, text=white, font=\scriptsize\bfseries, rounded corners=1mm,
    minimum width=0.7cm, anchor=north west},
}

% ---------------------------------------------------------------------
%  Diapositiva separadora de bloque
%  #1 número  #2 título  #3 duración  #4 color
% ---------------------------------------------------------------------
\newcommand{\bloque}[4]{%
\section{#2}
\begin{frame}[plain]
\begin{tikzpicture}[remember picture, overlay]
  \fill[#4] (current page.north west) rectangle
        ([xshift=5.2cm]current page.south west);
  \node[anchor=north west, text=white, font=\bfseries\large]
    at ([xshift=0.8cm,yshift=-2.2cm]current page.north west) {BLOQUE};
  \node[anchor=north west, text=white, font=\bfseries\fontsize{72}{72}\selectfont]
    at ([xshift=0.65cm,yshift=-2.9cm]current page.north west) {#1};
  \node[anchor=north west, text=iubnavy, font=\bfseries\LARGE,
        text width=9.4cm, align=left]
    at ([xshift=6.2cm,yshift=-2.8cm]current page.north west)
    {\hyphenpenalty=10000\exhyphenpenalty=10000 #2};
  \node[anchor=north west, fill=iubnavy, text=iubyellow, rounded corners=1.5mm,
        font=\small\bfseries, inner sep=2mm]
    at ([xshift=6.2cm,yshift=-6.0cm]current page.north west) {Duración: #3};
  \fill[iubnavy] ([xshift=5.2cm]current page.south west) rectangle
        ([yshift=0.35cm]current page.south east);
  \fill[iubyellow] ([xshift=5.2cm,yshift=0.35cm]current page.south west)
        rectangle ([yshift=0.42cm]current page.south east);
  \node[anchor=north east, text=iubnavy, font=\small\bfseries]
    at ([xshift=-0.6cm,yshift=-0.5cm]current page.north east) {IUB · IEL04};
\end{tikzpicture}
\end{frame}}

% =====================================================================
\begin{document}
% =====================================================================

% ---------------------------------------------------------------------
%  PORTADA
% ---------------------------------------------------------------------
\begin{frame}[plain]
\begin{tikzpicture}[remember picture, overlay]
  % Panel institucional
  \fill[iubnavy] (current page.north west) rectangle
        ([xshift=5.6cm]current page.south west);
  \node[anchor=north west, text=iubyellow, font=\bfseries\fontsize{54}{54}\selectfont]
    at ([xshift=0.7cm,yshift=-1.0cm]current page.north west) {IUB};
  \node[anchor=north west, text=white, font=\footnotesize, text width=4.3cm, align=left]
    at ([xshift=0.75cm,yshift=-3.0cm]current page.north west)
    {Institución Universitaria de Barranquilla};
  \draw[iubyellow, line width=1.2pt]
    ([xshift=0.75cm,yshift=-4.25cm]current page.north west) --
    ([xshift=2.6cm,yshift=-4.25cm]current page.north west);
  \node[anchor=north west, text=iubyellow, font=\scriptsize\bfseries,
        text width=4.3cm, align=left]
    at ([xshift=0.75cm,yshift=-4.5cm]current page.north west)
    {Tecnología en Gestión de Sistemas Eléctricos};
  \node[anchor=south west, text=white, font=\scriptsize, text width=4.3cm, align=left]
    at ([xshift=0.75cm,yshift=0.9cm]current page.south west)
    {Ing. Sergio Martinez Castilla\\[1pt]\textcolor{iubyellow}{\tiny smartinezc@unibarranquilla.edu.co}};
  % Bloque de título
  \node[anchor=north west, fill=iubblue, text=white, rounded corners=1.5mm,
        font=\scriptsize\bfseries, inner sep=2mm]
    at ([xshift=6.4cm,yshift=-1.0cm]current page.north west)
    {IEL04 \textperiodcentered{} Semana 11 de 14 \textperiodcentered{} Corte evaluativo 3};
  \node[anchor=north west, text=iubnavy, font=\small, text width=9.0cm, align=left] (a)
    at ([xshift=6.4cm,yshift=-1.9cm]current page.north west)
    {Circuitos Eléctricos I};
  \node[anchor=north west, text=iubnavy, font=\bfseries\fontsize{15}{18}\selectfont,
        text width=9.0cm, align=left] (t)
    at ([yshift=-0.35cm]a.south west)
    {\hyphenpenalty=10000 El circuito RLC en serie: reactancias opuestas, ley de voltajes de Kirchhoff y resonancia en serie};
  % Franjas de la paleta complementaria
  \fill[iubblue]   ([xshift=6.4cm,yshift=1.1cm]current page.south west)
    rectangle ++(1.6cm,0.18cm);
  \fill[iuborange] ([xshift=8.1cm,yshift=1.1cm]current page.south west)
    rectangle ++(1.6cm,0.18cm);
  \fill[iubgreen]  ([xshift=9.8cm,yshift=1.1cm]current page.south west)
    rectangle ++(1.6cm,0.18cm);
  \node[anchor=north west, text=iubnavy, font=\scriptsize]
    at ([xshift=6.4cm,yshift=0.95cm]current page.south west)
    {Sesión teórico-práctica \textperiodcentered{} 240 min};
\end{tikzpicture}
\end{frame}

% ---------------------------------------------------------------------
\begin{frame}{Punto de partida de la sesión}
\begin{columns}[T,onlytextwidth]
  \column{0.34\textwidth}
\begin{caja}[height=5.9cm]{iubblue}{Continuidad}
\scriptsize \begin{itemize}\setlength{\itemsep}{2pt}
  \item Semana 10: evaluación
  \item \textbf{Semana 11: Circuito RLC en serie: tipos e identificación de resistores, inductores y condensadores, análisis y Ley de Voltajes de Kirchhoff}
  \item Semana 12: Circuito RLC en paralelo
\end{itemize}
\end{caja}
  \column{0.34\textwidth}
\begin{caja}[height=5.9cm]{iuborange}{Prerrequisitos}
\scriptsize \begin{itemize}\setlength{\itemsep}{2pt}
  \item Calcular la reactancia capacitiva y la inductiva a una frecuencia dada
  \item Aplicar la ley de voltajes de Kirchhoff con fasores en circuitos RC
  \item Calcular la frecuencia de resonancia de un circuito LC (semana 3)
  \item Calcular potencia activa, reactiva y aparente y factor de potencia
\end{itemize}
\end{caja}
  \column{0.29\textwidth}
\begin{caja}[height=5.9cm]{iubgreen}{Competencia}
\scriptsize \textbf{UC2.} Coordinar actividades asociadas a sistemas eléctricos utilizando fuentes convencionales y no convencionales de energía.
\end{caja}
\end{columns}
\end{frame}

% ---------------------------------------------------------------------
\begin{frame}{Resultados de aprendizaje}
\begin{columns}[c,onlytextwidth]
  \column{0.45\textwidth}
\begin{caja}{iubnavy}{IEL04-RA4}
\footnotesize Calcular un circuito resistivo, inductivo, capacitivo y RLC, sea en serie o paralelo.
\end{caja}
\vspace{1mm}
\begin{cajaplana}{iubblue}
\scriptsize \textbf{CE1.} Identificar los tipos de resistores, inductores y capacitores.\\[2pt]
\textbf{CE2.} Realizar mediciones en un circuito resistivo, inductivo y capacitivo.\\[2pt]
\textbf{CE3.} Realizar mediciones en un circuito RLC.
\end{cajaplana}
  \column{0.52\textwidth}
\centering
\begin{adjustbox}{max width=\linewidth, max totalheight=6.1cm}
\begin{tikzpicture}
  \node[box, fill=iubblue, text width=4.9cm, align=left, font=\tiny, anchor=west] (r0) at (1.25,0.00) {Identificar el resistor, la bobina y el condensador de un circuito RLC y verificar sus valores, incluida la resistencia de devanado.};
  \node[tag=iubnavy, font=\tiny\bfseries, minimum width=1.15cm, anchor=east] at (1.1,0.00) {Aplicar};
  \node[box, fill=iuborange, text width=4.9cm, align=left, font=\tiny, anchor=west] (r1) at (1.25,-1.45) {Calcular la impedancia, el ángulo de fase y el carácter inductivo o capacitivo de un circuito RLC en serie a una frecuencia dada.};
  \node[tag=iubnavy, font=\tiny\bfseries, minimum width=1.15cm, anchor=east] at (1.1,-1.45) {Aplicar};
  \node[box, fill=iubgreen, text width=4.9cm, align=left, font=\tiny, anchor=west] (r2) at (1.25,-2.90) {Analizar las tensiones de un circuito RLC en serie con la ley de voltajes de Kirchhoff y explicar la resonancia en serie.};
  \node[tag=iubnavy, font=\tiny\bfseries, minimum width=1.15cm, anchor=east] at (1.1,-2.90) {Analizar};
  \node[box, fill=iubblue, text width=4.9cm, align=left, font=\tiny, anchor=west] (r3) at (1.25,-4.35) {Medir las tensiones en el resistor, la bobina y el condensador de un RLC en serie a varias frecuencias y verificar la resonancia.};
  \node[tag=iubnavy, font=\tiny\bfseries, minimum width=1.15cm, anchor=east] at (1.1,-4.35) {Evaluar};
  \draw[flecha, draw=iubnavy!40] (r0.south) -- (r1.north);
  \draw[flecha, draw=iubnavy!40] (r1.south) -- (r2.north);
  \draw[flecha, draw=iubnavy!40] (r2.south) -- (r3.north);
\end{tikzpicture}
\end{adjustbox}
\end{columns}
\end{frame}

% ---------------------------------------------------------------------
\begin{frame}{Ruta de la sesión \textperiodcentered{} 240 minutos}
\centering
\begin{tikzpicture}[x=1cm,y=1cm]
  \node[circle, fill=iubblue, text=white, font=\scriptsize\bfseries, minimum size=0.6cm, inner sep=0] at (0,0.00) {1};
  \node[anchor=west, font=\scriptsize, text width=7.6cm] at (0.45,0.00) {Del RC y el RL al RLC: reactancias de efectos opuestos y la resonancia};
  \fill[iubblue] (8.3,-0.20) rectangle ++(1.84,0.4);
  \node[anchor=west, font=\scriptsize\bfseries] at (10.24,0.00) {20 min};
  \node[circle, fill=iuborange, text=white, font=\scriptsize\bfseries, minimum size=0.6cm, inner sep=0] at (0,-0.76) {2};
  \node[anchor=west, font=\scriptsize, text width=7.6cm] at (0.45,-0.76) {Identificación y verificación de los tres componentes y elección de valores};
  \fill[iuborange] (8.3,-0.96) rectangle ++(1.84,0.4);
  \node[anchor=west, font=\scriptsize\bfseries] at (10.24,-0.76) {20 min};
  \node[circle, fill=iubgreen, text=white, font=\scriptsize\bfseries, minimum size=0.6cm, inner sep=0] at (0,-1.51) {3};
  \node[anchor=west, font=\scriptsize, text width=7.6cm] at (0.45,-1.51) {Z = R + j(XL \ensuremath{-} XC), reactancia total y carácter inductivo o capacitivo};
  \fill[iubgreen] (8.3,-1.71) rectangle ++(3.22,0.4);
  \node[anchor=west, font=\scriptsize\bfseries] at (11.62,-1.51) {35 min};
  \node[circle, fill=iubblue, text=white, font=\scriptsize\bfseries, minimum size=0.6cm, inner sep=0] at (0,-2.27) {4};
  \node[anchor=west, font=\scriptsize, text width=7.6cm] at (0.45,-2.27) {VS = VR + j(VL \ensuremath{-} VC): tensiones reactivas en oposición y mayores que la fuente};
  \fill[iubblue] (8.3,-2.47) rectangle ++(3.22,0.4);
  \node[anchor=west, font=\scriptsize\bfseries] at (11.62,-2.27) {35 min};
  \node[circle, fill=iuborange, text=white, font=\scriptsize\bfseries, minimum size=0.6cm, inner sep=0] at (0,-3.03) {5};
  \node[anchor=west, font=\scriptsize, text width=7.6cm] at (0.45,-3.03) {fr, Z = R, curva de corriente, Q = XL/R y ancho de banda};
  \fill[iuborange] (8.3,-3.23) rectangle ++(3.68,0.4);
  \node[anchor=west, font=\scriptsize\bfseries] at (12.08,-3.03) {40 min};
  \node[circle, fill=iubgreen, text=white, font=\scriptsize\bfseries, minimum size=0.6cm, inner sep=0] at (0,-3.79) {6};
  \node[anchor=west, font=\scriptsize, text width=7.6cm] at (0.45,-3.79) {Potencias reactivas que se compensan, FP unitario en resonancia};
  \fill[iubgreen] (8.3,-3.99) rectangle ++(2.30,0.4);
  \node[anchor=west, font=\scriptsize\bfseries] at (10.70,-3.79) {25 min};
  \node[circle, fill=iubblue, text=white, font=\scriptsize\bfseries, minimum size=0.6cm, inner sep=0] at (0,-4.54) {7};
  \node[anchor=west, font=\scriptsize, text width=7.6cm] at (0.45,-4.54) {Con el kit: 100 \ensuremath{\Omega}, bobina de 10.4 mH y condensador de 97.8 nF a 3 kHz, fr y 8 kHz};
  \fill[iubblue] (8.3,-4.74) rectangle ++(4.60,0.4);
  \node[anchor=west, font=\scriptsize\bfseries] at (13.00,-4.54) {50 min};
  \node[circle, fill=iuborange, text=white, font=\scriptsize\bfseries, minimum size=0.6cm, inner sep=0] at (0,-5.30) {8};
  \node[anchor=west, font=\scriptsize, text width=7.6cm] at (0.45,-5.30) {Síntesis, cierre actitudinal y bibliografía};
  \fill[iuborange] (8.3,-5.50) rectangle ++(1.38,0.4);
  \node[anchor=west, font=\scriptsize\bfseries] at (9.78,-5.30) {15 min};
  \draw[iubnavy, line width=0.6pt] (8.3,0.45) -- (8.3,-5.70);
\end{tikzpicture}
\end{frame}

% =====================================================================
\bloque{01}{Del RC y el RL al RLC: reactancias de efectos opuestos y la resonancia}{20 min}{iubblue}
% =====================================================================

% ---------------------------------------------------------------------
\begin{frame}{{\normalsize Tres elementos en serie: cuando las reactancias se oponen}}
\begin{columns}[c,onlytextwidth]
  \column{0.57\textwidth}
\small
\begin{itemize}\setlength{\itemsep}{4pt}
  \item Hasta la semana 9, cada circuito de corriente alterna tuvo un solo elemento reactivo: un condensador que adelantaba la corriente o una bobina que la retrasaba.
  \item En un receptor de radio, un circuito RLC en serie selecciona una emisora entre muchas porque a una sola frecuencia su impedancia es mínima.
\end{itemize}
  \column{0.4\textwidth}
\begin{cardO}{Nota pedagógica}
\footnotesize Conviene tener a mano ese resultado para compararlo con el de esta sesión.
\end{cardO}
\end{columns}
\end{frame}

% ---------------------------------------------------------------------
\begin{frame}{{\normalsize Del análisis del circuito RLC en serie a su medición}}
\centering
\begin{adjustbox}{max width=\linewidth, max totalheight=6.1cm}
\begin{tikzpicture}
  \node[bG, align=center, font=\scriptsize, text width=3.51cm] (n0) at (1.88,4.78) {R, L y C en serie};
  \node[bB, align=center, font=\scriptsize, text width=2.36cm] (n1) at (6.88,4.78) {XL y XC opuestas};
  \node[bB, align=center, font=\scriptsize, text width=2.80cm] (n2) at (11.54,4.78) {Z = R + j(XL \ensuremath{-} XC)};
  \node[bB, align=center, font=\scriptsize, text width=2.58cm] (n3) at (11.54,1.32) {LVK: VL y VC en oposición};
  \node[bO, align=center, font=\scriptsize, text width=3.03cm] (n4) at (6.88,1.32) {Resonancia: XL = XC};
  \node[bG, align=center, font=\scriptsize, text width=3.51cm] (n5) at (1.88,1.32) {Medir antes, en y después de fr};
  \draw[flecha] (n0) -- (n1);
  \draw[flecha] (n1) -- (n2);
  \draw[flecha] (n2) -- (n3);
  \draw[flecha] (n3) -- (n4);
  \draw[flecha] (n4) -- (n5);
\end{tikzpicture}
\end{adjustbox}

\vspace{1mm}{\scriptsize Del análisis del circuito RLC en serie a su medición\par}
\end{frame}

% =====================================================================
\bloque{02}{Identificación y verificación de los tres componentes y elección de valores}{20 min}{iuborange}
% =====================================================================

% ---------------------------------------------------------------------
\begin{frame}{{\normalsize Resistores, inductores y condensadores del circuito RLC}}
\begin{columns}[c,onlytextwidth]
  \column{0.57\textwidth}
\small
\begin{itemize}\setlength{\itemsep}{4pt}
  \item Para el caso de esta sesión se reutilizan dos componentes ya verificados y se agrega uno nuevo.
  \item La elección de un resistor de solo 100 \ensuremath{\Omega} no es casual.
  \item La tensión de trabajo del condensador también se comprueba antes de montar.
\end{itemize}
  \column{0.4\textwidth}
\begin{cardO}{Advertencia de seguridad}
\footnotesize En un circuito RLC en serie cerca de la resonancia, la tensión en el condensador y en la bobina puede ser mucho mayor que la de la fuente.
\end{cardO}
\end{columns}
\end{frame}

% ---------------------------------------------------------------------
\begin{frame}{{\normalsize Tipos de componentes del circuito RLC, qué se identifica}}
\centering\scriptsize
\renewcommand{\arraystretch}{1.35}
\rowcolors{2}{iubnavy!6}{white}
\begin{adjustbox}{max width=\linewidth, max totalheight=5.6cm}
\begin{tabularx}{\textwidth}{>{\raggedright\arraybackslash}X>{\raggedright\arraybackslash}X>{\raggedright\arraybackslash}X>{\raggedright\arraybackslash}X>{\raggedright\arraybackslash}X}
  \rowcolor{iubnavy}
  \hd{Componente} & \hd{Tipos principales} & \hd{Qué se lee en el cuerpo} & \hd{Qué se verifica} & \hd{Parámetro crítico en el RLC}\\
  Resistor & película de carbón, película metálica, bobinado, chip & código de colores o de tres dígitos & resistencia con ohmímetro & fija la corriente en resonancia y la agudeza\\
  Inductor & núcleo de aire, de hierro, de ferrita; fijo o variable & valor en µH o mH, código de colores & inductancia con LCR, RW con ohmímetro & L fija fr; RW se suma a R\\
  Condensador & mica, cerámico, película, electrolítico & valor, tolerancia, tensión de trabajo & capacitancia con LCR o capacímetro & C fija fr; su tensión nominal\\
\end{tabularx}
\end{adjustbox}
\vspace{2mm}
{\scriptsize Tipos de componentes del circuito RLC, qué se identifica en cada uno y con qué se verifica\par}
\end{frame}

% =====================================================================
\bloque{03}{Z = R + j(XL \ensuremath{-} XC), reactancia total y carácter inductivo o capacitivo}{35 min}{iubgreen}
% =====================================================================

% ---------------------------------------------------------------------
\begin{frame}{Impedancia del circuito RLC en serie}
\begin{columns}[c,onlytextwidth]
  \column{0.57\textwidth}
\small
\begin{itemize}\setlength{\itemsep}{4pt}
  \item En un circuito RLC en serie, la misma corriente atraviesa los tres elementos.
  \item A la izquierda, la impedancia sigue de cerca a $X_C$; a la derecha, a $X_L$.
  \item En la práctica, la resistencia $R$ de la fórmula no es solo la del resistor.
\end{itemize}
  \column{0.4\textwidth}
\begin{caja}{iubblue}{Ecuación clave}
\footnotesize \centering\small \begin{adjustbox}{max width=\linewidth}$\displaystyle \begin{gathered}\mathbf{Z}= R + jX_L - jX_C \\ = \sqrt{R^{2} + (X_L - X_C)^{2}}\ \angle \tan^{-1}\!\left(\frac{X_L - X_C}{R}\right)\end{gathered}$\end{adjustbox}
\end{caja}
\end{columns}
\end{frame}

% ---------------------------------------------------------------------
\begin{frame}{{\normalsize Reactancias XL y XC y magnitud de la impedancia de un RLC}}
\begin{columns}[c,onlytextwidth]
  \column{0.62\textwidth}
\centering
\begin{tikzpicture}
  \begin{semilogxaxis}[width=8.4cm, height=5.7cm, xmin=1000, xmax=20000, ymin=0, ymax=1683.3,
      xlabel={Frecuencia f (Hz)}, ylabel={Ohmios (\ensuremath{\Omega})},
      label style={font=\scriptsize, text=iubnavy}, tick label style={font=\tiny, text=iubnavy},
      axis line style={iubnavy}, grid=major, grid style={iubnavy!12},
      legend style={font=\tiny, fill=white}, legend pos=north east]
    \addplot[line width=1.4pt, iubblue] coordinates {(1000,62.832) (1040.9,65.401) (1083.5,68.075) (1127.8,70.859) (1173.9,73.756) (1221.9,76.772) (1271.8,79.912) (1323.8,83.179) (1378,86.581) (1434.3,90.121) (1493,93.806) (1554,97.642) (1617.6,101.63) (1683.7,105.79) (1752.6,110.12) (1824.2,114.62) (1898.8,119.31) (1976.5,124.18) (2057.3,129.26) (2141.4,134.55) (2229,140.05) (2320.1,145.78) (2415,151.74) (2513.7,157.94) (2616.5,164.4) (2723.5,171.12) (2834.9,178.12) (2950.8,185.4) (3071.4,192.98) (3197,200.88) (3327.8,209.09) (3463.8,217.64) (3605.5,226.54) (3752.9,235.8) (3906.4,245.44) (4066.1,255.48) (4232.4,265.93) (4405.4,276.8) (4585.6,288.12) (4773.1,299.9) (4968.2,312.16) (5171.4,324.93) (5382.9,338.22) (5603,352.05) (5832.1,366.44) (6070.6,381.42) (6318.8,397.02) (6577.2,413.26) (6846.1,430.15) (7126.1,447.74) (7417.4,466.05) (7720.7,485.11) (8036.5,504.95) (8365.1,525.59) (8707.1,547.08) (9063.2,569.46) (9433.8,592.74) (9819.5,616.98) (10221,642.21) (10639,668.47) (11074,695.8) (11527,724.25) (11998,753.87) (12489,784.69) (12999,816.78) (13531,850.18) (14084,884.94) (14660,921.13) (15260,958.79) (15884,998) (16533,1038.8) (17209,1081.3) (17913,1125.5) (18645,1171.5) (19408,1219.4) (20000,1256.6)};
    \addlegendentry{XL}
    \addplot[line width=1.4pt, iuborange] coordinates {(1000,1591.5) (1040.9,1529) (1083.5,1469) (1127.8,1411.3) (1173.9,1355.8) (1221.9,1302.6) (1271.8,1251.4) (1323.8,1202.2) (1378,1155) (1434.3,1109.6) (1493,1066) (1554,1024.2) (1617.6,983.92) (1683.7,945.27) (1752.6,908.13) (1824.2,872.46) (1898.8,838.18) (1976.5,805.26) (2057.3,773.62) (2141.4,743.23) (2229,714.03) (2320.1,685.98) (2415,659.03) (2513.7,633.15) (2616.5,608.27) (2723.5,584.38) (2834.9,561.42) (2950.8,539.37) (3071.4,518.18) (3197,497.82) (3327.8,478.26) (3463.8,459.48) (3605.5,441.43) (3752.9,424.08) (3906.4,407.42) (4066.1,391.42) (4232.4,376.04) (4405.4,361.27) (4585.6,347.08) (4773.1,333.44) (4968.2,320.34) (5171.4,307.76) (5382.9,295.67) (5603,284.05) (5832.1,272.9) (6070.6,262.18) (6318.8,251.88) (6577.2,241.98) (6846.1,232.47) (7126.1,223.34) (7417.4,214.57) (7720.7,206.14) (8036.5,198.04) (8365.1,190.26) (8707.1,182.79) (9063.2,175.61) (9433.8,168.71) (9819.5,162.08) (10221,155.71) (10639,149.6) (11074,143.72) (11527,138.07) (11998,132.65) (12489,127.44) (12999,122.43) (13531,117.62) (14084,113) (14660,108.56) (15260,104.3) (15884,100.2) (16533,96.264) (17209,92.482) (17913,88.849) (18645,85.359) (19408,82.006) (20000,79.577)};
    \addlegendentry{XC}
    \addplot[line width=1.4pt, iubgreen] coordinates {(1000,1532) (1040.9,1467) (1083.5,1404.4) (1127.8,1344.1) (1173.9,1286) (1221.9,1229.9) (1271.8,1175.7) (1323.8,1123.5) (1378,1073.1) (1434.3,1024.4) (1493,977.35) (1554,931.89) (1617.6,887.93) (1683.7,845.41) (1752.6,804.26) (1824.2,764.41) (1898.8,725.8) (1976.5,688.37) (2057.3,652.07) (2141.4,616.84) (2229,582.63) (2320.1,549.38) (2415,517.06) (2513.7,485.61) (2616.5,455) (2723.5,425.18) (2834.9,396.13) (2950.8,367.82) (3071.4,340.22) (3197,313.33) (3327.8,287.15) (3463.8,261.7) (3605.5,237.02) (3752.9,213.19) (3906.4,190.36) (4066.1,168.76) (4232.4,148.75) (4405.4,130.9) (4585.6,116.09) (4773.1,105.48) (4968.2,100.33) (5171.4,101.46) (5382.9,108.67) (5603,120.92) (5832.1,136.93) (6070.6,155.63) (6318.8,176.26) (6577.2,198.33) (6846.1,221.53) (7126.1,245.67) (7417.4,270.64) (7720.7,296.35) (8036.5,322.78) (8365.1,349.92) (8707.1,377.77) (9063.2,406.35) (9433.8,435.66) (9819.5,465.76) (10221,496.66) (10639,528.42) (11074,561.06) (11527,594.65) (11998,629.22) (12489,664.82) (12999,701.51) (13531,739.35) (14084,778.39) (14660,818.7) (15260,860.33) (15884,903.35) (16533,947.83) (17209,993.85) (17913,1041.5) (18645,1090.8) (19408,1141.8) (20000,1181.3)};
    \addlegendentry{Z}
    \addplot[line width=1pt, dashed, iubnavy!70] coordinates {(5033,0) (5033,1683.3)};
    \addlegendentry{fr = 5.03 kHz}
  \end{semilogxaxis}
\end{tikzpicture}
  \column{0.35\textwidth}
\begin{caja}{iubblue}{Lectura de la gráfica}
\footnotesize La gráfica muestra por qué la resonancia es un punto especial: la curva de impedancia tiene forma de V, con su mínimo exactamente donde se cruzan las curvas de $X_L$ y $X_C$.
\end{caja}
\end{columns}
\end{frame}

% =====================================================================
\bloque{04}{VS = VR + j(VL \ensuremath{-} VC): tensiones reactivas en oposición y mayores que la fuente}{35 min}{iubblue}
% =====================================================================

% ---------------------------------------------------------------------
\begin{frame}{{\normalsize Ley de voltajes de Kirchhoff en el circuito RLC en serie}}
\begin{columns}[c,onlytextwidth]
  \column{0.57\textwidth}
\small
\begin{itemize}\setlength{\itemsep}{4pt}
  \item Con la corriente común como referencia, la tensión del resistor está a 0°, la de la bobina a +90° y la del condensador a \ensuremath{-}90°.
  \item No viola ninguna ley: cada una es grande, pero como están en oposición, su efecto neto es pequeño.
  \item La fila de la resonancia es la más reveladora.
\end{itemize}
  \column{0.4\textwidth}
\begin{caja}{iubblue}{Ecuación clave}
\footnotesize \centering\small \begin{adjustbox}{max width=\linewidth}$\displaystyle \begin{gathered}\mathbf{V}_S = V_R + j(V_L - V_C) \\ V_S = \sqrt{V_R^{2} + (V_L - V_C)^{2}} \\ \theta = \tan^{-1}\!\left(\frac{V_L - V_C}{V_R}\right)\end{gathered}$\end{adjustbox}
\end{caja}
\end{columns}
\end{frame}

% ---------------------------------------------------------------------
\begin{frame}{Tensiones de un RLC en serie}
\centering\tiny
\renewcommand{\arraystretch}{1.35}
\rowcolors{2}{iubnavy!6}{white}
\begin{adjustbox}{max width=\linewidth, max totalheight=5.6cm}
\begin{tabularx}{\textwidth}{>{\raggedright\arraybackslash}X>{\raggedright\arraybackslash}X>{\raggedright\arraybackslash}X>{\raggedright\arraybackslash}X>{\raggedright\arraybackslash}X>{\raggedright\arraybackslash}X>{\raggedright\arraybackslash}X>{\raggedright\arraybackslash}X}
  \rowcolor{iubnavy}
  \hd{f (Hz)} & \hd{Z (\ensuremath{\Omega})} & \hd{I (mA)} & \hd{VR (V)} & \hd{VL (V)} & \hd{VC (V)} & \hd{\ensuremath{\surd}(VR² + (VL \ensuremath{-} VC)²) (V)} & \hd{Carácter}\\
  2000 & 677.5 & 2.95 & 0.30 & 0.37 & 2.35 & 2.00 & capacitivo\\
  4000 & 177.4 & 11.3 & 1.13 & 2.83 & 4.49 & 2.00 & capacitivo\\
  5033 (fr) & 100.0 & 20.0 & 2.00 & 6.32 & 6.32 & 2.00 & resistivo\\
  10 000 & 479.6 & 4.17 & 0.42 & 2.62 & 0.66 & 2.00 & inductivo\\
\end{tabularx}
\end{adjustbox}
\vspace{2mm}
{\scriptsize Tensiones de un RLC en serie (R = 100 \ensuremath{\Omega}, L = 10 mH, C = 0.1 µF, VS = 2 V eficaces) a varias frecuencias\par}
\end{frame}

% =====================================================================
\bloque{05}{fr, Z = R, curva de corriente, Q = XL/R y ancho de banda}{40 min}{iuborange}
% =====================================================================

% ---------------------------------------------------------------------
\begin{frame}{Resonancia en serie}
\begin{columns}[c,onlytextwidth]
  \column{0.57\textwidth}
\small
\begin{itemize}\setlength{\itemsep}{4pt}
  \item En el circuito en serie, la resonancia ocurre cuando $X_L = X_C$.
  \item Con $L = 10\ \text{mH}$, $C = 0.1\ \mu\text{F}$ y $V_S = 2\ \text{V}$: $f_r = 5.03\ \text{kHz}$ y $X_L = 316\ \Omega$.
\end{itemize}
  \column{0.4\textwidth}
\begin{caja}{iubblue}{Ecuación clave}
\footnotesize \centering\small \begin{adjustbox}{max width=\linewidth}$\displaystyle \begin{gathered}f_r = \frac{1}{2\pi\sqrt{LC}} \\ Z_r = R \\ I_{max} = \frac{V_S}{R} \\ Q = \frac{X_L}{R} \\ AB = \frac{f_r}{Q}\end{gathered}$\end{adjustbox}
\end{caja}
\end{columns}
\end{frame}

% ---------------------------------------------------------------------
\begin{frame}{Corriente de un RLC en serie}
\begin{columns}[c,onlytextwidth]
  \column{0.62\textwidth}
\centering
\begin{tikzpicture}
  \begin{semilogxaxis}[width=8.4cm, height=5.7cm, xmin=1000, xmax=25000, ymin=0.17065, ymax=20.949,
      xlabel={Frecuencia f (Hz)}, ylabel={Corriente I (mA)},
      label style={font=\scriptsize, text=iubnavy}, tick label style={font=\tiny, text=iubnavy},
      axis line style={iubnavy}, grid=major, grid style={iubnavy!12},
      legend style={font=\tiny, fill=white}, legend pos=north east]
    \addplot[line width=1.4pt, iubblue] coordinates {(1000,1.3055) (1044,1.3677) (1089.9,1.4334) (1137.9,1.5027) (1188,1.576) (1240.2,1.6536) (1294.8,1.7359) (1351.8,1.8232) (1411.3,1.916) (1473.4,2.0147) (1538.2,2.1201) (1605.9,2.2327) (1676.6,2.3533) (1750.3,2.4828) (1827.3,2.6222) (1907.8,2.7726) (1991.7,2.9355) (2079.3,3.1124) (2170.8,3.3054) (2266.4,3.5167) (2366.1,3.7492) (2470.2,4.0063) (2578.9,4.2922) (2692.4,4.612) (2810.8,4.9722) (2934.5,5.381) (3063.7,5.8489) (3198.5,6.3891) (3339.2,7.019) (3486.1,7.7613) (3639.5,8.6451) (3799.7,9.7078) (3966.9,10.994) (4141.4,12.55) (4323.7,14.396) (4513.9,16.464) (4712.5,18.464) (4919.9,19.797) (5136.4,19.836) (5362.4,18.561) (5598.4,16.579) (5844.7,14.505) (6101.9,12.644) (6370.4,11.072) (6650.7,9.7721) (6943.4,8.6983) (7248.9,7.8057) (7567.8,7.0565) (7900.9,6.4211) (8248.5,5.8764) (8611.5,5.405) (8990.4,4.9933) (9386,4.6306) (9799,4.3088) (10230,4.0212) (10680,3.7627) (11150,3.5289) (11641,3.3165) (12153,3.1225) (12688,2.9448) (13246,2.7812) (13829,2.6301) (14438,2.4902) (15073,2.3602) (15736,2.2391) (16429,2.1261) (17151,2.0203) (17906,1.9212) (18694,1.8281) (19517,1.7405) (20375,1.658) (21272,1.5802) (22208,1.5066) (23185,1.4371) (24205,1.3712) (25000,1.3241)};
    \addlegendentry{R = 100 \ensuremath{\Omega} (Q = 3.16)}
    \addplot[line width=1.4pt, iuborange] coordinates {(1000,1.2838) (1044,1.3428) (1089.9,1.4048) (1137.9,1.4699) (1188,1.5383) (1240.2,1.6102) (1294.8,1.6858) (1351.8,1.7654) (1411.3,1.8493) (1473.4,1.9376) (1538.2,2.0308) (1605.9,2.1291) (1676.6,2.2329) (1750.3,2.3426) (1827.3,2.4586) (1907.8,2.5812) (1991.7,2.7111) (2079.3,2.8487) (2170.8,2.9944) (2266.4,3.1488) (2366.1,3.3124) (2470.2,3.4857) (2578.9,3.6691) (2692.4,3.8629) (2810.8,4.0673) (2934.5,4.2821) (3063.7,4.5069) (3198.5,4.7406) (3339.2,4.9815) (3486.1,5.227) (3639.5,5.4734) (3799.7,5.7156) (3966.9,5.9472) (4141.4,6.1605) (4323.7,6.3468) (4513.9,6.4972) (4712.5,6.6034) (4919.9,6.659) (5136.4,6.6605) (5362.4,6.6078) (5598.4,6.5042) (5844.7,6.3561) (6101.9,6.1715) (6370.4,5.9594) (6650.7,5.7286) (6943.4,5.4868) (7248.9,5.2405) (7567.8,4.9949) (7900.9,4.7536) (8248.5,4.5195) (8611.5,4.2942) (8990.4,4.0788) (9386,3.8739) (9799,3.6795) (10230,3.4955) (10680,3.3217) (11150,3.1575) (11641,3.0026) (12153,2.8564) (12688,2.7185) (13246,2.5882) (13829,2.4651) (14438,2.3488) (15073,2.2388) (15736,2.1347) (16429,2.0361) (17151,1.9426) (17906,1.854) (18694,1.7699) (19517,1.6901) (20375,1.6142) (21272,1.5421) (22208,1.4735) (23185,1.4083) (24205,1.3461) (25000,1.3015)};
    \addlegendentry{R = 300 \ensuremath{\Omega} (Q = 1.05)}
    \addplot[line width=1pt, dashed, iubnavy!70] coordinates {(5033,0.17065) (5033,20.949)};
    \addlegendentry{fr = 5.03 kHz}
    \addplot[line width=1pt, dashed, iubnavy!70] coordinates {(1000,14.14) (25000,14.14)};
    \addlegendentry{0.707 · Imax (R = 100 \ensuremath{\Omega})}
  \end{semilogxaxis}
\end{tikzpicture}
  \column{0.35\textwidth}
\begin{caja}{iubblue}{Lectura de la gráfica}
\footnotesize Las dos curvas tienen el máximo en la misma frecuencia, porque $f_r$ solo depende de $L$ y $C$; la resistencia decide la altura y la anchura del pico.
\end{caja}
\end{columns}
\end{frame}

% =====================================================================
\bloque{06}{Potencias reactivas que se compensan, FP unitario en resonancia}{25 min}{iubgreen}
% =====================================================================

% ---------------------------------------------------------------------
\begin{frame}{{\normalsize Potencia y factor de potencia en el circuito RLC en serie}}
\begin{columns}[c,onlytextwidth]
  \column{0.57\textwidth}
\small
\begin{itemize}\setlength{\itemsep}{4pt}
  \item La columna de la resonancia resume lo más importante de la sesión en términos de energía.
  \item Esta compensación es la base de la \textbf{corrección del factor de potencia} en las instalaciones eléctricas.
\end{itemize}
  \column{0.4\textwidth}
\begin{caja}{iubblue}{Ecuación clave}
\footnotesize \centering\small \begin{adjustbox}{max width=\linewidth}$\displaystyle \begin{gathered}P = I^{2}R \\ Q = Q_L - Q_C = I^{2}(X_L - X_C) \\ S = V_S I = \sqrt{P^{2} + Q^{2}} \\ FP = \cos\theta = \frac{P}{S}\end{gathered}$\end{adjustbox}
\end{caja}
\end{columns}
\end{frame}

% ---------------------------------------------------------------------
\begin{frame}{Potencias de un RLC en serie}
\centering\scriptsize
\renewcommand{\arraystretch}{1.35}
\rowcolors{2}{iubnavy!6}{white}
\begin{adjustbox}{max width=\linewidth, max totalheight=5.6cm}
\begin{tabularx}{\textwidth}{>{\raggedright\arraybackslash}X>{\raggedright\arraybackslash}X>{\raggedright\arraybackslash}X>{\raggedright\arraybackslash}X}
  \rowcolor{iubnavy}
  \hd{Magnitud} & \hd{4 kHz} & \hd{5.03 kHz (fr)} & \hd{10 kHz}\\
  I (mA) & 11.3 & 20.0 & 4.17\\
  P = I²R (mW) & 12.7 & 40.0 & 1.74\\
  QL = I²XL (mVAR) & 31.9 & 126.5 & 10.9\\
  QC = I²XC (mVAR) & 50.5 & 126.5 & 2.77\\
  Q = QL \ensuremath{-} QC (mVAR) & \ensuremath{-}18.6 & 0 & 8.16\\
  S = VS·I (mVA) & 22.5 & 40.0 & 8.34\\
  FP = P/S & 0.564 en adelanto & 1.000 & 0.208 en atraso\\
\end{tabularx}
\end{adjustbox}
\vspace{2mm}
{\scriptsize Potencias de un RLC en serie (R = 100 \ensuremath{\Omega}, L = 10 mH, C = 0.1 µF, VS = 2 V eficaces)\par}
\end{frame}

% =====================================================================
\bloque{07}{Con el kit: 100 \ensuremath{\Omega}, bobina de 10.4 mH y condensador de 97.8 nF a 3 kHz, fr y 8 kHz}{50 min}{iubblue}
% =====================================================================

% ---------------------------------------------------------------------
\begin{frame}{{\normalsize Medición de un circuito RLC en serie antes, en y después}}
\begin{columns}[c,onlytextwidth]
  \column{0.57\textwidth}
\small
\begin{itemize}\setlength{\itemsep}{4pt}
  \item El generador entrega una onda sinusoidal ajustada a 2.00 V eficaces en los bornes del circuito, lo que elimina el efecto de su resistencia interna.
  \item La ley de voltajes de Kirchhoff queda entonces como $V_S = V_R + I R_W = 1.60 + 0.40 = 2.00\ \text{V}$: toda la tensión de la fuente se reparte entre las dos resistencias del lazo.
  \item En las otras dos frecuencias la verificación es la de la unidad sobre la LVK, con el ángulo real de la bobina.
\end{itemize}
  \column{0.4\textwidth}
\begin{caja}{iubblue}{Ecuación clave}
\footnotesize \centering\small \begin{adjustbox}{max width=\linewidth}$\displaystyle \begin{gathered}f_r= \frac{1}{2\pi\sqrt{(10.4 \times 10^{-3})(97.8 \times 10^{-9})}} \\ \approx 4.99\ \text{kHz} \\ Q = \frac{X_L}{R} = \frac{326.1}{124.1} \approx 2.63\end{gathered}$\end{adjustbox}
\end{caja}
\end{columns}
\end{frame}

% ---------------------------------------------------------------------
\begin{frame}{RLC en serie con los componentes del kit}
\centering\tiny
\renewcommand{\arraystretch}{1.35}
\rowcolors{2}{iubnavy!6}{white}
\begin{adjustbox}{max width=\linewidth, max totalheight=5.6cm}
\begin{tabularx}{\textwidth}{>{\raggedright\arraybackslash}X>{\raggedright\arraybackslash}X>{\raggedright\arraybackslash}X>{\raggedright\arraybackslash}X>{\raggedright\arraybackslash}X>{\raggedright\arraybackslash}X>{\raggedright\arraybackslash}X}
  \rowcolor{iubnavy}
  \hd{Magnitud} & \hd{3 kHz calc.} & \hd{3 kHz med.} & \hd{fr calc. (4.99 kHz)} & \hd{fr med. (4.97 kHz)} & \hd{8 kHz calc.} & \hd{8 kHz med.}\\
  Z (\ensuremath{\Omega}) & 368 \ensuremath{\angle} \ensuremath{-}70.3° & — & 124.1 \ensuremath{\angle} 0° & — & 343 \ensuremath{\angle} 68.8° & —\\
  I (mA) & 5.43 & 5.43 & 16.1 & 16.0 & 5.84 & 5.83\\
  VR (V) & 0.54 & 0.54 & 1.60 & 1.59 & 0.58 & 0.58\\
  Vbob (V) & 1.07 & 1.08 & 5.27 & 5.25 & 3.06 & 3.04\\
  VC (V) & 2.95 & 2.95 & 5.26 & 5.24 & 1.19 & 1.19\\
  Carácter & capacitivo & — & resistivo & — & inductivo & —\\
\end{tabularx}
\end{adjustbox}
\vspace{2mm}
{\scriptsize RLC en serie con los componentes del kit (VS = 2.00 V eficaces): cálculo con valores medidos y medición (lecturas de ejemplo)\par}
\end{frame}

% =====================================================================
\bloque{08}{Síntesis, cierre actitudinal y bibliografía}{15 min}{iuborange}
% =====================================================================

% ---------------------------------------------------------------------
\begin{frame}{Síntesis de la sesión}
\begin{columns}[T,onlytextwidth]
  \column{0.32\textwidth}
\begin{cardB}{Resistores, inductores y condensadores}
\scriptsize Identificación y verificación de los tres componentes y elección de valores
\end{cardB}
  \column{0.32\textwidth}
\begin{cardO}{Impedancia del circuito RLC en serie}
\scriptsize Z = R + j(XL \ensuremath{-} XC), reactancia total y carácter inductivo o capacitivo
\end{cardO}
  \column{0.32\textwidth}
\begin{cardG}{Ley de voltajes de Kirchhoff}
\scriptsize VS = VR + j(VL \ensuremath{-} VC): tensiones reactivas en oposición y mayores que la fuente
\end{cardG}
\end{columns}
\vspace{2mm}
\begin{columns}[T,onlytextwidth]
  \column{0.32\textwidth}
\begin{cardB}{Resonancia en serie}
\scriptsize fr, Z = R, curva de corriente, Q = XL/R y ancho de banda
\end{cardB}
  \column{0.32\textwidth}
\begin{cardO}{Potencia y factor de potencia}
\scriptsize Potencias reactivas que se compensan, FP unitario en resonancia
\end{cardO}
  \column{0.32\textwidth}
\begin{cardG}{Medición de un circuito RLC en serie}
\scriptsize Con el kit: 100 \ensuremath{\Omega}, bobina de 10.4 mH y condensador de 97.8 nF a 3 kHz, fr y 8 kHz
\end{cardG}
\end{columns}
\end{frame}

% ---------------------------------------------------------------------
\begin{frame}{Reflexión para llevar}
\centering
\begin{tikzpicture}
  \node[fill=iubnavy, text=white, rounded corners=6pt, text width=11.5cm, inner sep=12pt, align=center, font=\normalsize] (c) at (0,0) {\itshape «Revisar en equipo la tensión nominal del condensador antes de buscar la resonancia es responsabilidad compartida: en un RLC en serie, la tensión que aparece en sus bornes no la decide la fuente, sino el factor de calidad del circuito.»};
  \node[fill=iubyellow, text=iubnavy, rounded corners=3pt, font=\scriptsize\bfseries, inner sep=4pt, anchor=south west] at ([yshift=-4pt]c.north west) {Componente actitudinal};
\end{tikzpicture}
\end{frame}

% ---------------------------------------------------------------------
\begin{frame}{Referencias bibliográficas}
\small
\begin{itemize}\setlength{\itemsep}{8pt}
  \item[{$[1]$}] T. L. Floyd, Principios de circuitos eléctricos, 8.\textordfeminine{} ed. México: Pearson Educación, 2007.
  \item[{$[2]$}] R. L. Boylestad, Introducción al análisis de circuitos, 10.\textordfeminine{} ed. México: Pearson Educación, 2004.
  \item[{$[3]$}] M. F. P. Deorsola y P. Morcelle del Valle, Circuitos eléctricos. Parte 1, 1.\textordfeminine{} ed. La Plata: Editorial de la Universidad de La Plata, 2017.
\end{itemize}
\end{frame}

% ---------------------------------------------------------------------
%  CIERRE
% ---------------------------------------------------------------------
\begin{frame}[plain]
\begin{tikzpicture}[remember picture, overlay]
  \fill[iubnavy] (current page.north west) rectangle (current page.south east);
  \node[anchor=north west, text=iubyellow, font=\bfseries\fontsize{40}{44}\selectfont]
    at ([xshift=1.2cm,yshift=-1.6cm]current page.north west) {\textexclamdown{}Gracias!};
  \node[anchor=north west, text=white, font=\normalsize, text width=14cm, align=left]
    at ([xshift=1.2cm,yshift=-3.4cm]current page.north west)
    {IEL04 \textperiodcentered{} Circuitos Eléctricos I};
  \node[anchor=north west, fill=iubblue, text=white, rounded corners=1.5mm,
        font=\small\bfseries, inner sep=2.2mm, text width=11.5cm]
    at ([xshift=1.2cm,yshift=-4.4cm]current page.north west)
    {Próxima sesión \textperiodcentered{} Semana 12: Circuito RLC en paralelo: análisis y Ley de Corrientes de Kirchhoff};
  \node[anchor=south west, text=iubyellow, font=\small, align=left]
    at ([xshift=1.2cm,yshift=0.9cm]current page.south west)
    {Ing. Sergio Martinez Castilla \quad · \quad smartinezc@unibarranquilla.edu.co};
  \fill[iubblue]   ([xshift=1.2cm,yshift=0.55cm]current page.south west) rectangle ++(1.6cm,0.15cm);
  \fill[iuborange] ([xshift=2.9cm,yshift=0.55cm]current page.south west) rectangle ++(1.6cm,0.15cm);
  \fill[iubgreen]  ([xshift=4.6cm,yshift=0.55cm]current page.south west) rectangle ++(1.6cm,0.15cm);
\end{tikzpicture}
\end{frame}

\end{document}
